% Auto-generated complete chapter from 01_source_of_truth_profile.md
\chapter{01 — Source of Truth Profile}
\label{complete:root_002}

\begin{sourcemeta}
\textbf{Fuente:} \href{file:///E:/Laboral/01_source_of_truth_profile.md}{01\_source\_of\_truth\_profile.md}\\
\textbf{Seccion:} root\\
\textbf{Estado:} canonical\\
\textbf{Rol:} Modulo principal\\
\textbf{Lineas originales:} 929\\
\textbf{Ultima modificacion:} 2026-06-11 11:18\\
\textbf{Deduplicacion conservadora:} 0 bloques repetidos omitidos; 0 caracteres repetidos no reimpresos.
\end{sourcemeta}

\section*{Resumen de orientacion}
\textbf{Perfil:} Alexander Woodcock Salomón \textbf{Fecha de corte operativa:} 2026-06-11 \textbf{Estado del documento:} Fuente de verdad operativa. No es CV, LinkedIn, portfolio ni estrategia final. \textbf{Modo usado para esta sección:} análisis interno de archivos + contexto actualizado del usuario. \textbf{No requiere Deep Research ni modo agente.} \textbf{Regla crítica:} todo documento posterior debe respetar esta fuente de verdad. Si una sección posterior contradice este archivo, debe justificarlo c

\section*{Contenido completo depurado}
\addcontentsline{toc}{section}{Contenido completo depurado}


\subsection{Perfil maestro verificable para reconstrucción de estrategia laboral internacional}




\textbf{Perfil:} Alexander Woodcock Salomón  

\textbf{Fecha de corte operativa:} 2026-06-11  

\textbf{Estado del documento:} Fuente de verdad operativa. No es CV, LinkedIn, portfolio ni estrategia final.  

\textbf{Modo usado para esta sección:} análisis interno de archivos + contexto actualizado del usuario. \textbf{No requiere Deep Research ni modo agente.}  

\textbf{Regla crítica:} todo documento posterior debe respetar esta fuente de verdad. Si una sección posterior contradice este archivo, debe justificarlo con nueva evidencia o corregirse.




\medskip\hrule\medskip




\subsection{0. Propósito}




Este archivo fija los hechos, límites, matices y claims defendibles sobre Alexander Woodcock Salomón antes de reconstruir el resto de la investigación laboral.




Su función es evitar que los documentos posteriores mezclen:




\begin{itemize}
\item información vigente con cronogramas antiguos;
\item hechos demostrables con aspiraciones;
\item habilidades defendibles con habilidades en desarrollo;
\item rutas migratorias futuras con derechos laborales actuales;
\item cursos adquiridos con certificaciones formales;
\item proyectos funcionales con proyectos incompletos o no auditados;
\item estrategia principal con rutas secundarias o de contingencia.
\end{itemize}




Este documento \textbf{no decide todavía la estrategia laboral final}. Define el perímetro de verdad desde el cual se redactarán las secciones posteriores.




\medskip\hrule\medskip




\subsection{1. Fuentes usadas y estatus}




\begin{center}
\scriptsize
\begin{longtable}{>{\raggedright\arraybackslash}p{0.225\linewidth} >{\raggedright\arraybackslash}p{0.225\linewidth} >{\raggedright\arraybackslash}p{0.225\linewidth} >{\raggedright\arraybackslash}p{0.225\linewidth}}
\toprule
Fuente & Tipo & Uso en este archivo & Estatus \\
\midrule
\endfirsthead
\toprule
Fuente & Tipo & Uso en este archivo & Estatus \\
\midrule
\endhead
\texttt{Pegado text(12).txt} & Prompt original largo & Perfil, objetivos, proyectos, restricciones, entregables esperados & Vigente como fuente primaria de intención \\
\texttt{deep-research-report (8).md} & Informe anterior & Material a auditar: mercado, CV, LinkedIn, GitHub, portfolio, scripts, movilidad, salarios & Útil pero insuficiente; no debe tratarse como verdad cerrada \\
\texttt{alexander\_woodcock\_context.md} & Contexto maestro antiguo & Identidad, restricciones, hardware, herramientas, runway antiguo, stack, perfil psicológico, planes & Parcialmente histórico; debe revalidarse \\
\texttt{plan\_audit.md} & Auditoría antigua de tesis/portfolio & Criterios de optimización, documentación de profiling, riesgo operativo & Parcialmente obsoleto; tesis ya no está en desarrollo \\
\texttt{master\_schedule.md} & Cronograma antiguo & Evidencia de prioridades, proyectos y secuenciación previa & Obsoleto como calendario \\
\texttt{ticktick\_schedule.md} & Agenda granular antigua & Evidencia de tareas planeadas, cursos y prioridades técnicas & Obsoleto como calendario \\
\texttt{job\_search\_strategy.md} & Estrategia antigua & Hipótesis “Trojan Horse”: Junior Technical Artist / QA Automation / Tools & Hipótesis útil; requiere revalidación \\
\texttt{portfolio\_strategy.md} & Estrategia antigua de portfolio & Optimization Doctor, Toolsmith, Math Wizard, Mobile Hero, Human Portrait & Útil como arquitectura conceptual; requiere priorización \\
\texttt{estrategia\_idiomas\_mercado\_laboral\_technical\_artist.md} & Investigación de idiomas/mercados & Base para estrategia de inglés, alemán, portugués, francés, UE & Útil; requiere actualización en secciones externas \\
\texttt{Courses 2025.xlsx} & Inventario de cursos adquiridos & Insumo para matriz de formación/ROI & Debe analizarse en la sección 06; no son certificaciones \\
\texttt{00\_master\_audit\_blueprint.md} & Auditoría maestra ya generada & Marco estructural de reconstrucción & Vigente \\
\bottomrule
\end{longtable}
\end{center}




\medskip\hrule\medskip




\subsection{2. Estado maestro congelado}




\subsubsection{2.1 Identidad}




\begin{center}
\scriptsize
\begin{longtable}{>{\raggedright\arraybackslash}p{0.455\linewidth} >{\raggedright\arraybackslash}p{0.455\linewidth}}
\toprule
Campo & Valor operativo \\
\midrule
\endfirsthead
\toprule
Campo & Valor operativo \\
\midrule
\endhead
Nombre completo & Alexander Woodcock Salomón \\
Ubicación & Colombia \\
Ciudad/región reportada & Pasto, Nariño, Colombia \\
Zona horaria & UTC-5 / Colombia \\
Idioma nativo & Español \\
Inglés & C1 autoevaluado; funcional para entrevistas técnicas si se sostiene en práctica \\
Alemán & Ninguno / no profesional \\
Portugués & Estratégicamente relevante, no competencia profesional confirmada \\
Perfil actual deseado & Technical Artist / Real-Time 3D Developer / Unity WebGL / CAD-to-realtime / Tools \\
\bottomrule
\end{longtable}
\end{center}




\subsubsection{2.2 Estado académico}




\begin{center}
\scriptsize
\begin{longtable}{>{\raggedright\arraybackslash}p{0.455\linewidth} >{\raggedright\arraybackslash}p{0.455\linewidth}}
\toprule
Elemento & Valor operativo \\
\midrule
\endfirsthead
\toprule
Elemento & Valor operativo \\
\midrule
\endhead
Carrera actual & Ingeniería Multimedia — UNAD \\
Estado de tesis & \textbf{Tesis completa, pendiente de sustentación} \\
Proyecto de tesis & TwinSight X500 / WebGL-Thesis-Proposal \\
Carrera previa & Ingeniería Electrónica — Universidad Nacional de Colombia, UNAL \\
Estado UNAL & Aproximadamente 8 semestres / cerca de 80\% del currículo, no título final confirmado \\
Forma internacional recomendada para UNAL & \texttt{Completed advanced coursework in Electronic Engineering — Universidad Nacional de Colombia} \\
Forma que debe evitarse & \texttt{Electronic Engineer}, \texttt{BSc Electronic Engineering}, \texttt{degree completed}, si el título no fue obtenido \\
\bottomrule
\end{longtable}
\end{center}




\subsubsection{2.3 Estado laboral}




\begin{center}
\scriptsize
\begin{longtable}{>{\raggedright\arraybackslash}p{0.455\linewidth} >{\raggedright\arraybackslash}p{0.455\linewidth}}
\toprule
Elemento & Valor operativo \\
\midrule
\endfirsthead
\toprule
Elemento & Valor operativo \\
\midrule
\endhead
Experiencia formal en Technical Art / Unity / Real-Time 3D & Baja / no fuerte todavía \\
Experiencia freelance & Sí, especialmente web/proyectos simples, pero no debe venderse como experiencia industrial senior \\
Experiencia no target & Call center / YMCA camps / experiencia en inglés y comunicación \\
Principal evidencia profesional actual & Proyectos: TwinSight X500, ARA Framework, Blender portrait, proyectos web/AI secundarios \\
Nivel profesional defendible & Entry-level fuerte / junior fuerte / early-career technical artist-real-time developer, según rol \\
Nivel que no debe afirmarse & Senior Technical Artist, Lead, Production Engineer, EU-based worker, industry-proven specialist \\
\bottomrule
\end{longtable}
\end{center}




\subsubsection{2.4 Estado legal y movilidad}




\begin{center}
\scriptsize
\begin{longtable}{>{\raggedright\arraybackslash}p{0.455\linewidth} >{\raggedright\arraybackslash}p{0.455\linewidth}}
\toprule
Elemento & Valor operativo \\
\midrule
\endfirsthead
\toprule
Elemento & Valor operativo \\
\midrule
\endhead
Nacionalidad actual relevante & Colombiana, salvo evidencia adicional no confirmada aquí \\
Autorización laboral UE actual & No confirmada. Operativamente: \textbf{no existe} \\
Ciudadanía/pasaporte portugués & Expectativa futura / proceso estratégico, no derecho actual \\
Ruta Alemania por pareja/matrimonio & Escenario posible futuro, condicionado a hechos no actuales \\
Prioridad de corto plazo & Remoto desde Colombia \\
Modelo laboral más realista inicial & Contractor/B2B, freelance internacional, remote-first, EOR si empresa lo ofrece \\
Claim prohibido & \texttt{Authorized to work in the EU} hasta que sea legalmente cierto \\
\bottomrule
\end{longtable}
\end{center}




\subsubsection{2.5 Estado financiero y operativo}




\begin{center}
\scriptsize
\begin{longtable}{>{\raggedright\arraybackslash}p{0.455\linewidth} >{\raggedright\arraybackslash}p{0.455\linewidth}}
\toprule
Elemento & Estado \\
\midrule
\endfirsthead
\toprule
Elemento & Estado \\
\midrule
\endhead
Runway financiero & Datos antiguos indican presión y bajo capital; debe actualizarse \\
Deadline financiero antiguo & Archivos antiguos hablan de May 2026; tratar como histórico hasta reconfirmación \\
OPEX & Vive con padres / bajo costo relativo según contexto anterior; debe reconfirmarse \\
Urgencia laboral & Alta pero no debe forzar mala estrategia ni sobreventa \\
Setup contractor & No confirmado: RUT, facturación electrónica, contador, Wise/Payoneer/Deel/Remote, ARL/EPS/pensión \\
\bottomrule
\end{longtable}
\end{center}




\textbf{Regla:} ninguna sección posterior debe hacer cálculos de supervivencia, runway o cadencia de aplicaciones sin reconfirmar datos financieros actuales.




\medskip\hrule\medskip




\subsection{3. Corrección de premisas desactualizadas}




Los archivos antiguos son útiles como contexto, pero no como estado actual. La siguiente tabla fija cómo deben interpretarse.




\begin{center}
\scriptsize
\begin{longtable}{>{\raggedright\arraybackslash}p{0.225\linewidth} >{\raggedright\arraybackslash}p{0.225\linewidth} >{\raggedright\arraybackslash}p{0.225\linewidth} >{\raggedright\arraybackslash}p{0.225\linewidth}}
\toprule
Tema en archivos antiguos & Estado anterior documentado & Estado operativo actual & Uso permitido \\
\midrule
\endfirsthead
\toprule
Tema en archivos antiguos & Estado anterior documentado & Estado operativo actual & Uso permitido \\
\midrule
\endhead
Retopología del dron & Pendiente / no iniciada & Ya no debe asumirse pendiente & Solo como historial de proceso \\
UVs / baking / Unity import & Planeados & No asumir pendientes sin nueva confirmación & Verificar si quedan assets o capturas por generar \\
Testing SUS / NASA-TLX & Planeado & Métricas ya reportadas en prompt & Confirmar versión final de datos antes de publicar \\
Entrega tesis Feb 6 / fases Dec-May & Cronograma viejo & Obsoleto & No usar como plan actual \\
Job hunt Mar 9 / Panic Protocol May 1 & Plan viejo & Obsoleto en fechas & Puede inspirar contingencia, no agenda vigente \\
Portfolio Bridge Feb-Mar & Plan viejo & Obsoleto & Reconvertir en backlog moderno \\
Cursos paralelos durante tesis & Plan viejo & Ya no aplica como restricción de tesis & Mantener como inventario de aprendizaje \\
Shaders creados / VFX reel planeado & Contexto útil & Requiere verificar estado real & No afirmar como portfolio publicado \\
\bottomrule
\end{longtable}
\end{center}




\textbf{Regla operativa:} toda mención a “tesis en desarrollo” debe reemplazarse por “tesis completa, pendiente de sustentación”, salvo cuando se describa retrospectivamente el proceso técnico.




\medskip\hrule\medskip




\subsection{4. Matriz de verdad: hechos, pendientes, supuestos, aspiraciones y no afirmables}




\subsubsection{4.1 Hechos defendibles actualmente}




\begin{center}
\scriptsize
\begin{longtable}{>{\raggedright\arraybackslash}p{0.305\linewidth} >{\raggedright\arraybackslash}p{0.305\linewidth} >{\raggedright\arraybackslash}p{0.305\linewidth}}
\toprule
Hecho & Nivel de confianza & Uso permitido \\
\midrule
\endfirsthead
\toprule
Hecho & Nivel de confianza & Uso permitido \\
\midrule
\endhead
Alexander estudia Ingeniería Multimedia en UNAD & Alto & CV, LinkedIn, portfolio \\
La tesis está completa y pendiente de sustentación & Alto según actualización del usuario & CV con wording prudente: \texttt{Thesis completed; defense pending} \\
TwinSight X500 / WebGL-Thesis-Proposal es el proyecto principal & Alto & Principal caso de estudio \\
El proyecto usa Unity WebGL / C\# / Blender / CAD-to-realtime & Alto según prompt/contexto & Posicionamiento, CV, LinkedIn \\
El proyecto transforma documentación/CAD de dron en visualización técnica 3D interactiva & Alto & Portfolio, entrevistas \\
Se reportan métricas: 6.5M+ triángulos a 95,617, SUS 91.88, NASA-TLX 8.69 vs 19.89, 12 participantes, 96 registros & Medio-Alto; requiere confirmar versión final antes de publicación & Bullets cuantificados si son finales y verificables \\
ARA Framework existe como prototipo de automatización de investigación con LLMs & Medio-Alto & Proyecto secundario, AI tooling \\
Hay cursos adquiridos sin certificación formal & Alto & Desarrollo profesional, no certificación \\
Inglés C1 es autoevaluado & Medio & \texttt{Professional working proficiency} si se sostiene oralmente; evitar sobreprometer si no se ha validado \\
La prioridad laboral es remoto desde Colombia & Alto & Estrategia principal \\
\bottomrule
\end{longtable}
\end{center}




\subsubsection{4.2 Datos pendientes de verificar}




\begin{center}
\scriptsize
\begin{longtable}{>{\raggedright\arraybackslash}p{0.305\linewidth} >{\raggedright\arraybackslash}p{0.305\linewidth} >{\raggedright\arraybackslash}p{0.305\linewidth}}
\toprule
Dato & Por qué importa & Dónde se verificará \\
\midrule
\endfirsthead
\toprule
Dato & Por qué importa & Dónde se verificará \\
\midrule
\endhead
URL pública funcional de TwinSight & Necesaria para portfolio y aplicaciones & Sección 08 / revisión GitHub-demo \\
Estado actual del repo \texttt{WebGL-Thesis-Proposal} & Determina limpieza, README y riesgo de assets & Sección 10 / GitHub audit \\
Estado de ArtStation & Define si sirve para aplicar o debe ocultarse temporalmente & Sección 07 / 10 \\
Estado actual de LinkedIn real & El informe anterior se basa en capturas/contexto parcial & Sección 10 \\
Estado de ARA Framework & Decide si publicarlo o mantenerlo privado & Sección 09 \\
Estado de \texttt{ai-news-aggregator} & Decide si vale pulirlo o descartarlo & Sección 09 \\
Estado de \texttt{FitApp-Free} & Probablemente no es prioritario; confirmar & Sección 09 \\
Estado del retrato Blender & Determina si puede publicarse como breakdown técnico & Sección 07 / 09 \\
Cursos completados vs solo adquiridos & Evita presentar acumulación como competencia real & Sección 06 \\
Acceso actual a Rebelway & Determina ROI real de cursos de VFX/Houdini & Sección 06 \\
Setup legal Colombia para contractor & Necesario antes de aceptar pagos internacionales & Sección 03 / 04 \\
Estado real de la ruta portuguesa & No debe inferirse & Sección 04 \\
Runway financiero actual & Define agresividad del plan & Sección 14 \\
Fecha real de sustentación & Define cadencia de portfolio y aplicaciones & Sección 14 \\
\bottomrule
\end{longtable}
\end{center}




\subsubsection{4.3 Supuestos estratégicos válidos pero no cerrados}




\begin{center}
\scriptsize
\begin{longtable}{>{\raggedright\arraybackslash}p{0.305\linewidth} >{\raggedright\arraybackslash}p{0.305\linewidth} >{\raggedright\arraybackslash}p{0.305\linewidth}}
\toprule
Supuesto & Estado & Riesgo \\
\midrule
\endfirsthead
\toprule
Supuesto & Estado & Riesgo \\
\midrule
\endhead
El mejor eje inicial es Technical Artist / Real-Time 3D Developer & Plausible & Debe validarse con mercado y roles reales \\
Technical Visualization / Digital Twin / Simulation puede tener mejor ROI que game tech art puro & Plausible & Requiere datos de salarios y demanda \\
USD 1,500/mes es alcanzable rápido & Plausible & Depende de mercado, portfolio y modalidad contractor \\
USD 3,000/mes en 3–6 meses es posible pero no garantizado & Hipótesis & Debe validarse con benchmarks LATAM/global remote \\
USD 6,000/mes en 12–24 meses es alcanzable si hay portafolio fuerte y experiencia acumulada & Hipótesis ambiciosa & Requiere ruta clara a mid-level o nicho mejor pagado \\
ARA puede abrir ruta secundaria AI tooling & Plausible & Puede diluir el perfil si se sobrepresenta \\
Rebelway/Houdini puede mejorar tech art/VFX a mediano plazo & Plausible & Riesgo de dispersión si se prioriza antes del portfolio principal \\
Alemán es inversión estratégica por DACH/UE & Plausible & Bajo ROI en 0–6 meses; requiere horizonte largo \\
\bottomrule
\end{longtable}
\end{center}




\subsubsection{4.4 Aspiraciones legítimas}




\begin{center}
\scriptsize
\begin{longtable}{>{\raggedright\arraybackslash}p{0.455\linewidth} >{\raggedright\arraybackslash}p{0.455\linewidth}}
\toprule
Aspiración & Cómo debe formularse \\
\midrule
\endfirsthead
\toprule
Aspiración & Cómo debe formularse \\
\midrule
\endhead
Trabajar remoto desde Colombia para empresa extranjera & Objetivo principal de corto plazo \\
Llegar a USD 3,000/mes & Meta realista condicionada, no promesa \\
Llegar a USD 6,000/mes & Meta de 12–24 meses, requiere evidencia y mejor mercado \\
Eventual movilidad UE & Ruta estratégica futura, no derecho actual \\
Aplicar a roles de videojuegos & Interés personal, pero no necesariamente mejor ROI inicial \\
Aplicar a simulación, digital twin, technical visualization & Ruta racional si paga mejor y encaja con TwinSight \\
Aprender alemán/portugués & Inversión mediano-largo plazo, no prioridad por encima de inglés+portfolio en corto plazo \\
\bottomrule
\end{longtable}
\end{center}




\subsubsection{4.5 Afirmaciones prohibidas o de alto riesgo}




\begin{center}
\scriptsize
\begin{longtable}{>{\raggedright\arraybackslash}p{0.455\linewidth} >{\raggedright\arraybackslash}p{0.455\linewidth}}
\toprule
Afirmación & Motivo \\
\midrule
\endfirsthead
\toprule
Afirmación & Motivo \\
\midrule
\endhead
\texttt{EU work authorized} & No existe autorización confirmada actualmente \\
\texttt{Portuguese citizen} & No está confirmado actualmente \\
\texttt{German resident} & No está confirmado actualmente \\
\texttt{Certified by Rebelway / CG Cookie / Alive} & Cursos sin certificado formal no son certificaciones \\
\texttt{Senior Technical Artist} & No hay experiencia industrial formal suficiente \\
\texttt{Production-ready AI research platform} para ARA & Riesgo de sobreventa \\
\texttt{Full-stack developer} como identidad principal & Diluye el posicionamiento 3D/Unity si no hay evidencia fuerte \\
\texttt{Industry-proven Unity developer} & No hay experiencia formal industrial fuerte aún \\
\texttt{Built everything manually without AI} & Falso si hubo asistencia IA \\
\texttt{AI built the project} & También falso si hubo dirección, integración, depuración y validación propia \\
\bottomrule
\end{longtable}
\end{center}




\medskip\hrule\medskip




\subsection{5. Posicionamiento operativo provisional}




\subsubsection{5.1 Posicionamiento base}




\begin{quote}
\textbf{Technical Artist \& Real-Time 3D Developer especializado en Unity, WebGL, C\#, Blender, CAD-to-realtime optimization, interactive technical visualization y Python/AI-assisted production tooling.}
\end{quote}




Este posicionamiento es provisional pero fuerte porque conecta:




\begin{itemize}
\item Unity/C\#;
\item WebGL/Web/mobile deployment;
\item Blender/asset optimization;
\item CAD-to-realtime;
\item visualización técnica;
\item experiencia de tesis;
\item base de ingeniería electrónica;
\item tooling/automation con Python e IA.
\end{itemize}




\subsubsection{5.2 Qué NO debe ser el posicionamiento principal}




\begin{center}
\scriptsize
\begin{longtable}{>{\raggedright\arraybackslash}p{0.455\linewidth} >{\raggedright\arraybackslash}p{0.455\linewidth}}
\toprule
Posicionamiento & Problema \\
\midrule
\endfirsthead
\toprule
Posicionamiento & Problema \\
\midrule
\endhead
Multimedia Designer & Demasiado genérico y mal pagado \\
3D Generalist & Compite contra artistas con reels más fuertes; no aprovecha ingeniería \\
Game Artist & Falta portfolio artístico específico de videojuegos \\
Full-stack Developer & Ruta saturada, desvía de TwinSight \\
AI Engineer & ARA no basta para competir contra perfiles ML/software fuertes \\
VFX Houdini Artist & Cursos disponibles, pero no evidencia profesional actual suficiente \\
UX/UI Designer & TwinSight tiene UI, pero no es la identidad principal \\
\bottomrule
\end{longtable}
\end{center}




\subsubsection{5.3 Narrativa primaria sugerida}




\begin{quote}
Perfil de ingeniería multimedia con base previa en electrónica, especializado en convertir modelos CAD y sistemas técnicos en experiencias 3D interactivas en tiempo real usando Unity, WebGL, C\#, Blender y optimización de assets.
\end{quote}




\subsubsection{5.4 Narrativa secundaria permitida}




\begin{quote}
Además del eje 3D/Unity, desarrolla herramientas Python e IA para automatizar investigación, documentación y flujos de producción.
\end{quote}




\subsubsection{5.5 Jerarquía de roles provisional}




\begin{center}
\scriptsize
\begin{longtable}{>{\raggedright\arraybackslash}p{0.455\linewidth} >{\raggedright\arraybackslash}p{0.455\linewidth}}
\toprule
Nivel & Roles \\
\midrule
\endfirsthead
\toprule
Nivel & Roles \\
\midrule
\endhead
Primarios & Unity Technical Artist, Real-Time 3D Developer, Interactive 3D Developer, Unity WebGL Developer, Technical Visualization Developer \\
Secundarios fuertes & CAD-to-Realtime Technical Artist, Digital Twin Visualization Developer, XR/Simulation Developer, Tools Developer for art/production pipelines \\
Secundarios tácticos & Python AI Tools Developer, Research Automation Developer, QA Automation technical entry route \\
Contingencia & Full-stack junior, frontend/React, bilingual support/BPO temporal, freelance Unity bug fixing \\
No prioritarios ahora & Pure 3D Artist, pure VFX Houdini Artist, pure game artist, pure data analyst \\
\bottomrule
\end{longtable}
\end{center}




\textbf{Nota:} esta jerarquía debe validarse en \texttt{02\_market\_role\_fit\_and\_positioning.md} con investigación de mercado actual.




\medskip\hrule\medskip




\subsection{6. Inventario maestro de proyectos}




\subsubsection{6.1 TwinSight X500 / WebGL-Thesis-Proposal}




\begin{center}
\scriptsize
\begin{longtable}{>{\raggedright\arraybackslash}p{0.455\linewidth} >{\raggedright\arraybackslash}p{0.455\linewidth}}
\toprule
Campo & Valor operativo \\
\midrule
\endfirsthead
\toprule
Campo & Valor operativo \\
\midrule
\endhead
Prioridad & Máxima \\
Tipo & Tesis / flagship portfolio project \\
Dominio & Unity WebGL, visualización técnica, CAD-to-realtime, technical UI \\
Tema & Visualizador web 3D interactivo para inspección/ensamblaje de dron Holybro X500 V2 \\
Claim central & “Seeing parts is not enough. Users need to understand spatial relationships.” \\
Stack reportado & Unity, C\#, WebGL, URP, UI Toolkit, Shader Graph, Blender, CAD/STEP pipeline \\
Funciones reportadas & selección de componentes, exploded view, clipping/cross-section, paneles técnicos, modos visuales, mobile/web performance \\
Métricas reportadas & 6.5M+ triángulos → 95,617; SUS 91.88; NASA-TLX 8.69 vs 19.89; 12 participantes; 96 records \\
Contribución personal & programación/integración Unity, interacción, UI, optimización, documentación, evaluación, dirección técnica con asistencia IA \\
Riesgos & assets de fabricante/CAD, estado público del repo, demo funcional, explicación de IA, métricas finales verificables \\
\bottomrule
\end{longtable}
\end{center}




\paragraph{Uso permitido}



\begin{itemize}
\item Proyecto principal de CV.
\item Primer caso de estudio del portfolio.
\item Primer Featured item en LinkedIn.
\item Primer pinned repo si está limpio.
\item Base para roles Unity WebGL, Interactive 3D, Technical Visualization, CAD-to-realtime, Digital Twin, Simulation.
\end{itemize}




\paragraph{Claims defendibles si se verifican}



\begin{itemize}
\item \texttt{Developed a Unity WebGL interactive 3D technical visualization prototype for drone assembly inspection.}
\item \texttt{Optimized CAD-derived assets for real-time web deployment.}
\item \texttt{Implemented interactive inspection features including component selection, exploded view and cross-section/clipping.}
\item \texttt{Evaluated usability and perceived workload using SUS, NASA-TLX Raw and Think-Aloud.}
\item \texttt{Reduced CAD-derived geometry from 6.5M+ triangles to 95,617 optimized triangles.}
\end{itemize}




\paragraph{Claims que requieren cuidado}



\begin{center}
\scriptsize
\begin{longtable}{>{\raggedright\arraybackslash}p{0.455\linewidth} >{\raggedright\arraybackslash}p{0.455\linewidth}}
\toprule
Claim & Condición \\
\midrule
\endfirsthead
\toprule
Claim & Condición \\
\midrule
\endhead
\texttt{Mobile-first} & Solo si la demo fue diseñada/probada realmente para móvil \\
\texttt{Thermal visualization} & Aclarar si es cualitativa/visual, no simulación física/FEA \\
\texttt{Digital twin} & Usar solo si se define como visualization/prototype, no gemelo digital industrial completo \\
\texttt{AI-assisted implementation} & Explicar sin debilitar credibilidad ni ocultar uso de IA \\
\texttt{Academic validation} & Usar si protocolos/resultados están completos y documentados \\
\bottomrule
\end{longtable}
\end{center}




\paragraph{Claims prohibidos}



\begin{itemize}
\item \texttt{Production-ready industrial digital twin}.
\item \texttt{FEA thermal simulation} si no hay análisis físico real.
\item \texttt{Commercial product deployed to clients} si no es cierto.
\item \texttt{Built all source code manually without AI}.
\item \texttt{Senior-level Unity architecture} si no se valida con revisión técnica.
\end{itemize}




\medskip\hrule\medskip




\subsubsection{6.2 ARA Framework}




\begin{center}
\scriptsize
\begin{longtable}{>{\raggedright\arraybackslash}p{0.455\linewidth} >{\raggedright\arraybackslash}p{0.455\linewidth}}
\toprule
Campo & Valor operativo \\
\midrule
\endfirsthead
\toprule
Campo & Valor operativo \\
\midrule
\endhead
Prioridad & Secundaria \\
Tipo & AI research automation prototype \\
Dominio & Python, LangGraph, LLM orchestration, multi-agent workflow, research automation \\
Estado reportado & Casi funcional / necesita pulido \\
Stack reportado & Python 3.12+, LangGraph, LangChain, Groq LLaMA 3.3-70B, Redis, Semantic Scholar, Playwright MCP, MarkItDown MCP, PDF processing \\
Arquitectura & 5 agentes: Niche Analyst, Literature Researcher, Technical Architect, Implementation Specialist, Content Synthesizer \\
Riesgo & Diluir posicionamiento principal; sobreventa como producto enterprise \\
\bottomrule
\end{longtable}
\end{center}




\paragraph{Uso permitido}



\begin{itemize}
\item Proyecto secundario para demostrar tooling mindset.
\item Ruta alternativa para Python/AI tools.
\item Soporte a narrativa de automatización de producción/investigación.
\item Puede estar en portfolio secundario si tiene README, demo y límites claros.
\end{itemize}




\paragraph{Claims defendibles}



\begin{itemize}
\item \texttt{Prototype for AI-assisted academic research automation.}
\item \texttt{Designed a multi-agent workflow using LangGraph and Python.}
\item \texttt{Integrated academic search, browser automation and markdown report generation components.}
\end{itemize}




\paragraph{Claims prohibidos o de alto riesgo}



\begin{itemize}
\item \texttt{Production-ready research automation platform}.
\item \texttt{Enterprise AI product}.
\item \texttt{Autonomous scientific research system}.
\item \texttt{Reliable academic citation engine} sin validación fuerte.
\end{itemize}




\medskip\hrule\medskip




\subsubsection{6.3 ai-news-aggregator}




\begin{center}
\scriptsize
\begin{longtable}{>{\raggedright\arraybackslash}p{0.455\linewidth} >{\raggedright\arraybackslash}p{0.455\linewidth}}
\toprule
Campo & Valor operativo \\
\midrule
\endfirsthead
\toprule
Campo & Valor operativo \\
\midrule
\endhead
Prioridad & Baja-media, según estado real \\
Tipo & Full-stack / AI news aggregation project \\
Stack reportado & FastAPI, React, PostgreSQL, Redis, Python backend \\
Estado reportado & Funcionó antes, necesita ajuste \\
Uso potencial & Ruta secundaria full-stack/AI \\
Riesgo & Distrae del posicionamiento Unity/3D si se sobredimensiona \\
\bottomrule
\end{longtable}
\end{center}




\paragraph{Uso permitido}



\begin{itemize}
\item Solo incluir si hay demo funcional, README, screenshots y arquitectura limpia.
\item Útil para roles Python/AI tooling o full-stack secundario.
\item No debe desplazar a TwinSight.
\end{itemize}




\paragraph{Reglas}



\begin{itemize}
\item Si no corre: mantener privado.
\item Si corre parcialmente: presentarlo como prototype.
\item Si no tiene relación con visualización/3D/tools: no poner en CV principal.
\end{itemize}




\medskip\hrule\medskip




\subsubsection{6.4 FitApp-Free}




\begin{center}
\scriptsize
\begin{longtable}{>{\raggedright\arraybackslash}p{0.455\linewidth} >{\raggedright\arraybackslash}p{0.455\linewidth}}
\toprule
Campo & Valor operativo \\
\midrule
\endfirsthead
\toprule
Campo & Valor operativo \\
\midrule
\endhead
Prioridad & Baja para estrategia Technical Artist \\
Tipo & Fitness/product/app logic project \\
Estado & En desarrollo \\
Uso potencial & Product systems thinking / full-stack contingency \\
Riesgo & Alta dispersión; poca relación con Unity/Technical Art \\
\bottomrule
\end{longtable}
\end{center}




\paragraph{Uso permitido}



\begin{itemize}
\item No incluir en CV principal para Technical Artist.
\item Podría incluirse en portfolio secundario solo si tiene demo y una narrativa de sistemas/product logic.
\item No dedicarle tiempo antes del portfolio TwinSight, LinkedIn, CV y aplicaciones.
\end{itemize}




\medskip\hrule\medskip




\subsubsection{6.5 Hyperrealistic Blender Portrait / CG Cookie Human}




\begin{center}
\scriptsize
\begin{longtable}{>{\raggedright\arraybackslash}p{0.455\linewidth} >{\raggedright\arraybackslash}p{0.455\linewidth}}
\toprule
Campo & Valor operativo \\
\midrule
\endfirsthead
\toprule
Campo & Valor operativo \\
\midrule
\endhead
Prioridad & Media \\
Tipo & 3D art / character pipeline / technical breakdown \\
Fuente de aprendizaje & CG Cookie HUMAN — completado según Excel/contexto \\
Valor & Demuestra Blender, topology, grooming, materials, lighting, end-to-end character workflow \\
Riesgo & Si se presenta solo como render, compite como character artist y diluye perfil técnico \\
\bottomrule
\end{longtable}
\end{center}




\paragraph{Uso correcto}



Presentarlo como:




\begin{quote}
\textbf{Technical Art / real-time adaptation breakdown}, no como “soy character artist puro”.
\end{quote}




Debe mostrar:




\begin{itemize}
\item topology;
\item UVs;
\item bake maps;
\item materials;
\item grooming/hair strategy;
\item lighting;
\item optimization considerations;
\item posibilidad de mobile/real-time conversion si se realiza.
\end{itemize}




\paragraph{Uso incorrecto}



\begin{itemize}
\item Publicar solo beauty render sin breakdown.
\item Usarlo como pieza principal del portfolio por encima de TwinSight.
\item Afirmar experiencia profesional de character artist si solo fue curso/proyecto personal.
\end{itemize}




\medskip\hrule\medskip




\subsubsection{6.6 Proyectos/repositories antiguos o menores}




\begin{center}
\scriptsize
\begin{longtable}{>{\raggedright\arraybackslash}p{0.305\linewidth} >{\raggedright\arraybackslash}p{0.305\linewidth} >{\raggedright\arraybackslash}p{0.305\linewidth}}
\toprule
Proyecto/repositorio & Estado operativo & Uso recomendado provisional \\
\midrule
\endfirsthead
\toprule
Proyecto/repositorio & Estado operativo & Uso recomendado provisional \\
\midrule
\endhead
\texttt{deus-ex-machina} & Viejo / GameMaker-style & Auditar. Probablemente no destacar salvo que tenga valor histórico o README contextual. \\
\texttt{skills-introduction-to-github} & Tutorial & Ocultar/no pinear/no mostrar a reclutadores. \\
\texttt{Chocolateria} & Irrelevante para estrategia actual & Ignorar salvo razón fuerte. \\
\texttt{Bosque de Sombras Manuales} & Irrelevante & Ignorar. \\
\texttt{alexander-fullstack-portfolio} & Test/prototipo & Auditar si sirve para portfolio site. \\
\texttt{alexander-3d-portfolio} & Test/prototipo & Auditar si sirve para portfolio site. \\
\texttt{delarge95.github.io} & Potencial portfolio host & Alta utilidad si se convierte en sitio principal. \\
\texttt{3DWeb} & Potencial web 3D & Auditar. \\
\texttt{job-search-strategy} & Documentación interna & No mostrar como repo profesional si contiene información personal. \\
\texttt{ludex\_framework} & Desconocido & Auditar antes de decidir. \\
\texttt{WoodSleiman} & Desconocido & Auditar antes de decidir. \\
\texttt{WebGL\_tesis} & Relacionado con tesis & Auditar duplicación con repo principal. \\
\bottomrule
\end{longtable}
\end{center}




\medskip\hrule\medskip




\subsection{7. Inventario maestro de habilidades}




\subsubsection{7.1 Habilidades primarias defendibles}




Estas habilidades pueden aparecer en la narrativa principal si están respaldadas por TwinSight, tesis, código o portfolio.




\begin{center}
\scriptsize
\begin{longtable}{>{\raggedright\arraybackslash}p{0.455\linewidth} >{\raggedright\arraybackslash}p{0.455\linewidth}}
\toprule
Categoría & Habilidades \\
\midrule
\endfirsthead
\toprule
Categoría & Habilidades \\
\midrule
\endhead
Engine / realtime & Unity, C\#, Unity WebGL, URP, UI Toolkit \\
3D / asset pipeline & Blender, retopology, UVs, baking, high-poly to low-poly, CAD-to-realtime optimization \\
Rendering / visual modes & Shader Graph, basic HLSL, real-time rendering concepts, x-ray/blueprint/solid/wireframe/ghosted visual modes if implemented \\
Technical visualization & interactive 3D UI, technical inspection, assembly visualization, cross-section/clipping \\
Performance & triangle reduction, mesh optimization, web/mobile performance considerations, profiling concepts \\
Documentation / evaluation & technical documentation, SUS, NASA-TLX Raw, Think-Aloud, academic validation \\
Tooling & Git/GitHub, Python, AI-assisted coding/productivity \\
Foundations & multimedia engineering, computer graphics fundamentals, math/physics/electronics background from UNAL \\
\bottomrule
\end{longtable}
\end{center}




\subsubsection{7.2 Habilidades secundarias que requieren validación por proyecto}




\begin{center}
\scriptsize
\begin{longtable}{>{\raggedright\arraybackslash}p{0.455\linewidth} >{\raggedright\arraybackslash}p{0.455\linewidth}}
\toprule
Habilidad & Evidencia requerida antes de afirmarla fuerte \\
\midrule
\endfirsthead
\toprule
Habilidad & Evidencia requerida antes de afirmarla fuerte \\
\midrule
\endhead
FastAPI & Demo running / repo / endpoint documentation \\
React & UI functional / deployed app / component structure \\
PostgreSQL & schema, migrations, queries, deployment proof \\
Redis & real usage, not only dependency \\
LangGraph & running graph, agent states, checkpointing evidence \\
LLM orchestration & demos, logs, architecture diagram \\
Multi-agent systems & stable agent workflow, reproducible output \\
Playwright/MCP & scripts or recorded automation \\
Semantic Scholar integrations & API calls, rate limit handling, output validation \\
Full-stack apps & deployed or locally reproducible app \\
GameMaker/GML & old project evidence; likely not central \\
Houdini/VFX & course outputs or reel, not just access \\
Unreal/Niagara & course outputs or demonstrable project \\
Three.js/WebGPU & project proof if added later \\
\bottomrule
\end{longtable}
\end{center}




\subsubsection{7.3 Habilidades aspiracionales o de mediano plazo}




\begin{center}
\scriptsize
\begin{longtable}{>{\raggedright\arraybackslash}p{0.455\linewidth} >{\raggedright\arraybackslash}p{0.455\linewidth}}
\toprule
Habilidad & Condición para mover a defendible \\
\midrule
\endfirsthead
\toprule
Habilidad & Condición para mover a defendible \\
\midrule
\endhead
Advanced HLSL & shader code examples or technical breakdown \\
WebGPU & prototype/demo \\
Three.js / React Three Fiber & web 3D project \\
XR foundations & Unity XR prototype, headset/mobile AR test \\
Houdini procedural tools & HDA/demo, not course only \\
Realtime VFX & flipbook, Niagara/Unity VFX demo, reel \\
Pipeline tools & Unity EditorWindow, Blender Python tool, batch processor \\
Technical writing & published case studies, README quality \\
German B1/B2 & exam or functional interview capability \\
Portuguese A2/B1 & practical use/learning status \\
\bottomrule
\end{longtable}
\end{center}




\subsubsection{7.4 Habilidades que no deben liderar el perfil ahora}




\begin{center}
\scriptsize
\begin{longtable}{>{\raggedright\arraybackslash}p{0.455\linewidth} >{\raggedright\arraybackslash}p{0.455\linewidth}}
\toprule
Habilidad / identidad & Motivo \\
\midrule
\endfirsthead
\toprule
Habilidad / identidad & Motivo \\
\midrule
\endhead
Data Analyst & No conecta directamente con TwinSight ni portfolio 3D actual \\
Pure Full-stack Developer & Mercado saturado; evidencia no principal \\
Pure AI Engineer / ML Engineer & Falta base formal/proyectos ML robustos \\
Pure Houdini FX Artist & Mucho acceso a cursos, poca evidencia profesional actual \\
Pure Character Artist & Retrato Blender es útil, pero no suficiente como eje \\
UX/UI Designer & UI Toolkit ayuda, pero no es el valor diferencial principal \\
\bottomrule
\end{longtable}
\end{center}




\medskip\hrule\medskip




\subsection{8. Formación, cursos y credenciales}




\subsubsection{8.1 Política general}




Los cursos adquiridos son un activo de aprendizaje y producción de portfolio. \textbf{No son certificaciones}, salvo que exista certificado verificable, institución formal y credencial emitida.




Uso correcto:




\begin{itemize}
\item \texttt{Professional development}.
\item \texttt{Self-directed training}.
\item \texttt{Coursework / self-study in [topic]}.
\item Mencionar cursos solo si están conectados a un output visible.
\end{itemize}




Uso incorrecto:




\begin{itemize}
\item \texttt{Certified Technical Artist}.
\item \texttt{Rebelway Certified}.
\item \texttt{CG Cookie Certified}.
\item Enumerar decenas de cursos en CV.
\item Presentar acceso a cursos como competencia profesional ya adquirida.
\end{itemize}




\subsubsection{8.2 Cursos y áreas detectadas en \texttt{Courses 2025.xlsx}}




El Excel contiene una lista extensa de cursos adquiridos, con columnas como:




\begin{itemize}
\item Nombre;
\item Área;
\item Software;
\item Academia;
\item Año;
\item Link;
\item Archivo;
\item Clave;
\item Enfoque;
\item Habilidades;
\item Weeks;
\item Video Hours;
\item Estado;
\item Notas.
\end{itemize}




Áreas detectadas:




\begin{center}
\scriptsize
\begin{longtable}{>{\raggedright\arraybackslash}p{0.455\linewidth} >{\raggedright\arraybackslash}p{0.455\linewidth}}
\toprule
Área & Uso potencial \\
\midrule
\endfirsthead
\toprule
Área & Uso potencial \\
\midrule
\endhead
Blender / 3D Art & Soporte a asset pipeline, modeling, topology, shading \\
Hard Surface & Útil para drones, props, CAD-to-realtime, game-ready assets \\
Character pipeline & Útil si se convierte en technical breakdown \\
Houdini / VFX & Mediano-largo plazo para proceduralism, FX, technical art \\
Unreal / Niagara & Secundario; útil si se abre ruta Unreal/VFX \\
Compositing / Nuke / DaVinci & Útil para demo reel, presentation polish, VFX; bajo impacto CV principal \\
Python scripting & Útil para tools developer, Blender/Houdini automation \\
Web 3D / Spline / Blender to Web & Útil si se conecta a Unity WebGL / interactive 3D \\
AI / Machine Learning & Útil si se conecta a ARA o tooling, no como ML engineer claim \\
\bottomrule
\end{longtable}
\end{center}




\subsubsection{8.3 Cursos específicos relevantes ya identificados}




\begin{center}
\scriptsize
\begin{longtable}{>{\raggedright\arraybackslash}p{0.305\linewidth} >{\raggedright\arraybackslash}p{0.305\linewidth} >{\raggedright\arraybackslash}p{0.305\linewidth}}
\toprule
Curso / grupo & Estado reportado & Valor operativo \\
\midrule
\endfirsthead
\toprule
Curso / grupo & Estado reportado & Valor operativo \\
\midrule
\endhead
CG Cookie HUMAN — Realistic Portrait Creation with Blender & Completado según Excel & Útil para Blender portrait / pipeline breakdown \\
Alive! — 3D Animation Mechanics \& Dynamic Motion & Curso adquirido/no certificado según prompt & Útil para fundamentos de animación, no credential fuerte \\
Blender Bros Hard Surface / Hard Ops / Box Cutter & Adquiridos según Excel & Útiles para hard-surface asset pipeline \\
Interactive Content Creation Unleashed: From Blender to Web & Adquirido & Potencial conexión con web 3D / thesis \\
CG Cookie scripting / Blender Python & Adquirido & Potencial para tools/pipeline \\
Houdini-Course.com & En progreso según Excel & Mediano-largo plazo; no eje actual sin output \\
Rebelway Houdini Fundamentals & Acceso detectado & Alto valor de skill, alto riesgo de dispersión \\
Rebelway Intro to UE & Acceso detectado & Útil si se decide Unreal/VFX; no prioridad sobre Unity actual \\
Rebelway Advanced Asset Creation & Acceso detectado & Potencial para procedural assets \\
Rebelway Stylized Realtime FX for Games & Acceso detectado & Potencial si se produce realtime FX reel \\
Rebelway Python for Houdini Artists & Acceso detectado & Potencial alto para Tools/VFX pipeline \\
Rebelway VEX for Houdini Artists & Acceso detectado & Mediano-largo plazo \\
Rebelway Intro to Machine Learning & Pausado según Excel & No priorizar antes de fundamentos/proyecto claro \\
\bottomrule
\end{longtable}
\end{center}




\subsubsection{8.4 Regla de ROI para cursos}




Un curso solo pasa a la estrategia principal si produce al menos uno de estos outputs:




\begin{enumerate}
\item portfolio case study;
\item demo reel segment;
\item shader/tool breakdown;
\item GitHub repo limpio;
\item ArtStation technical breakdown;
\item improvement visible en TwinSight;
\item evidence for a specific job title.
\end{enumerate}




Sin output, el curso se trata como consumo formativo, no como activo laboral.




\medskip\hrule\medskip




\subsection{9. Idiomas y movilidad: fuente de verdad operacional}




\subsubsection{9.1 Idiomas actuales}




\begin{center}
\scriptsize
\begin{longtable}{>{\raggedright\arraybackslash}p{0.305\linewidth} >{\raggedright\arraybackslash}p{0.305\linewidth} >{\raggedright\arraybackslash}p{0.305\linewidth}}
\toprule
Idioma & Estado operativo & Uso laboral \\
\midrule
\endfirsthead
\toprule
Idioma & Estado operativo & Uso laboral \\
\midrule
\endhead
Español & Nativo & LATAM, España, Colombia, networking local \\
Inglés & C1 autoevaluado & Principal idioma laboral internacional \\
Alemán & No profesional & No usar como ventaja laboral actual \\
Portugués & Relevante estratégicamente & No usar como competencia profesional actual salvo evidencia \\
Francés / otros & No confirmados & No relevantes actualmente \\
\bottomrule
\end{longtable}
\end{center}




\subsubsection{9.2 Regla de corto plazo}




En 0–6 meses, el orden correcto es:




\begin{enumerate}
\item defensa/sustentación;
\item portfolio TwinSight;
\item CV/LinkedIn/GitHub;
\item aplicaciones y networking;
\item inglés técnico/interview practice;
\item cursos solo si producen portfolio;
\item idiomas adicionales solo en intensidad baja.
\end{enumerate}




\subsubsection{9.3 Movilidad UE}




\begin{center}
\scriptsize
\begin{longtable}{>{\raggedright\arraybackslash}p{0.305\linewidth} >{\raggedright\arraybackslash}p{0.305\linewidth} >{\raggedright\arraybackslash}p{0.305\linewidth}}
\toprule
Ruta & Estado & Uso en estrategia \\
\midrule
\endfirsthead
\toprule
Ruta & Estado & Uso en estrategia \\
\midrule
\endhead
Portugal / pasaporte futuro & No confirmado como derecho actual & Plan mediano/largo plazo \\
Alemania por pareja/matrimonio & Condicionado a hechos futuros & Plan contingente, no base operativa \\
EU local jobs & No asumir elegibilidad actual & Aplicar solo si remote/contractor o sponsorship posible \\
Remote desde Colombia para EU company & Posible según empresa, contractor/EOR & Prioridad investigable \\
\bottomrule
\end{longtable}
\end{center}




\subsubsection{9.4 Claim permitido a reclutadores}




Versión prudente:




\begin{quote}
Based in Colombia and available for remote international contractor roles. Open to future EU mobility if legally viable; no current EU work authorization should be assumed.
\end{quote}




Versión que debe evitarse:




\begin{quote}
I am EU work authorized / I can work in Germany / I will have Portuguese citizenship soon.
\end{quote}




\medskip\hrule\medskip




\subsection{10. Restricciones operativas y riesgos personales}




\subsubsection{10.1 Riesgos reportados en contexto antiguo}




\begin{center}
\scriptsize
\begin{longtable}{>{\raggedright\arraybackslash}p{0.455\linewidth} >{\raggedright\arraybackslash}p{0.455\linewidth}}
\toprule
Riesgo & Implicación estratégica \\
\midrule
\endfirsthead
\toprule
Riesgo & Implicación estratégica \\
\midrule
\endhead
Analysis paralysis & Limitar opciones y producir entregables concretos \\
Shiny Object Syndrome & Evitar muchos cursos/proyectos paralelos \\
Perfeccionismo & Definir MVP de portfolio y deadlines \\
Bajo capital / presión financiera antigua & Reconfirmar runway y diseñar plan realista \\
Hardware con GTX 980 Ti & Evitar depender de UE5/RT/heavy rendering como eje inmediato \\
Tendencia a construir sistemas grandes & Priorizar empleabilidad antes que arquitectura perfecta \\
\bottomrule
\end{longtable}
\end{center}




\subsubsection{10.2 Reglas de mitigación}




\begin{center}
\scriptsize
\begin{longtable}{>{\raggedright\arraybackslash}p{0.455\linewidth} >{\raggedright\arraybackslash}p{0.455\linewidth}}
\toprule
Regla & Justificación \\
\midrule
\endfirsthead
\toprule
Regla & Justificación \\
\midrule
\endhead
TwinSight primero & Es el activo más alineado con el mercado \\
Portfolio MVP antes que portfolio perfecto & Empleabilidad > perfeccionismo \\
Máximo 2 proyectos visibles al inicio & Reduce dispersión narrativa \\
Cursos solo con output & Evita consumo infinito \\
No pivotar a AI/full-stack salvo decisión estratégica & Evita dilución \\
No usar EU como excusa para retrasar aplicaciones remotas & EU es mediano plazo \\
No construir 100 documentos antes de aplicar & Los documentos deben habilitar acción, no reemplazarla \\
\bottomrule
\end{longtable}
\end{center}




\medskip\hrule\medskip




\subsection{11. Mapa de evidencia requerida por activo}




\subsubsection{11.1 Evidencia mínima para TwinSight}




\begin{center}
\scriptsize
\begin{longtable}{>{\raggedright\arraybackslash}p{0.455\linewidth} >{\raggedright\arraybackslash}p{0.455\linewidth}}
\toprule
Evidencia & Prioridad \\
\midrule
\endfirsthead
\toprule
Evidencia & Prioridad \\
\midrule
\endhead
Demo video 60–120 s & Muy alta \\
Case study escrito & Muy alta \\
Screenshots UI principales & Muy alta \\
Before/after optimization table & Alta \\
GitHub README recruiter-friendly & Alta \\
Technical README / architecture diagram & Media-alta \\
SUS / NASA-TLX results visualization & Alta si los datos son finales \\
Explanation of AI-assisted development & Alta \\
Asset/legal note for CAD/manufacturer models & Alta \\
WebGL live demo & Muy alta si funciona estable \\
\bottomrule
\end{longtable}
\end{center}




\subsubsection{11.2 Evidencia mínima para ARA}




\begin{center}
\scriptsize
\begin{longtable}{>{\raggedright\arraybackslash}p{0.455\linewidth} >{\raggedright\arraybackslash}p{0.455\linewidth}}
\toprule
Evidencia & Prioridad \\
\midrule
\endfirsthead
\toprule
Evidencia & Prioridad \\
\midrule
\endhead
README con alcance y límites & Alta \\
Architecture diagram & Alta \\
Example input/output & Alta \\
Running demo or screen recording & Media-alta \\
Security/API key cleanup & Muy alta antes de publicar \\
Clear “prototype” wording & Muy alta \\
\bottomrule
\end{longtable}
\end{center}




\subsubsection{11.3 Evidencia mínima para Blender portrait}




\begin{center}
\scriptsize
\begin{longtable}{>{\raggedright\arraybackslash}p{0.455\linewidth} >{\raggedright\arraybackslash}p{0.455\linewidth}}
\toprule
Evidencia & Prioridad \\
\midrule
\endfirsthead
\toprule
Evidencia & Prioridad \\
\midrule
\endhead
Beauty render & Media \\
Wireframe/topology & Alta \\
UV layout & Alta \\
Material/shader breakdown & Alta \\
Grooming/lighting breakdown & Media-alta \\
Optimization/real-time adaptation notes & Alta si se usa para technical art \\
\bottomrule
\end{longtable}
\end{center}




\subsubsection{11.4 Evidencia mínima para LinkedIn/CV/GitHub}




\begin{center}
\scriptsize
\begin{longtable}{>{\raggedright\arraybackslash}p{0.455\linewidth} >{\raggedright\arraybackslash}p{0.455\linewidth}}
\toprule
Activo & Evidencia mínima \\
\midrule
\endfirsthead
\toprule
Activo & Evidencia mínima \\
\midrule
\endhead
CV & 1-page ATS version, 4 variants later \\
LinkedIn & headline, About, experience, projects, featured, skills \\
GitHub & pinned repos, README, screenshots, repo hygiene, no tutorial repos destacados \\
Portfolio & homepage, TwinSight case study, contact CTA, project links \\
\bottomrule
\end{longtable}
\end{center}




\medskip\hrule\medskip




\subsection{12. Narrative stack: qué historia debe contar el perfil}




\subsubsection{12.1 Historia principal}




Alexander debe presentarse como un perfil técnico-creativo que une ingeniería, 3D y software para construir experiencias interactivas de visualización técnica.




Estructura narrativa:




\begin{enumerate}
\item base previa en ingeniería electrónica;
\item transición a ingeniería multimedia;
\item especialización en Unity/WebGL/Blender;
\item tesis como prueba fuerte: TwinSight X500;
\item capacidad de convertir CAD/documentación técnica en visualización 3D interactiva;
\item capacidad secundaria de automatizar con Python/IA;
\item objetivo: roles remotos internacionales en real-time 3D, technical visualization, Unity, XR/simulation/digital twin.
\end{enumerate}




\subsubsection{12.2 Historia secundaria}




ARA y proyectos AI muestran que Alexander no es solo artista 3D ni solo Unity user, sino alguien con mentalidad de tooling, automation y research systems.




Debe formularse así:




\begin{quote}
I also build Python/AI-assisted tooling prototypes for research automation and production workflows.
\end{quote}




No debe formularse así:




\begin{quote}
I am an AI engineer specialized in autonomous agents.
\end{quote}




\subsubsection{12.3 Historia de transición académica}




Forma correcta:




\begin{quote}
I completed advanced coursework in Electronic Engineering at Universidad Nacional de Colombia before moving into Multimedia Engineering, where I focused on computer graphics, real-time 3D, interaction and technical visualization.
\end{quote}




Forma incorrecta:




\begin{quote}
I left engineering and became a designer.
\end{quote}




\subsubsection{12.4 Historia sobre IA asistida}




Forma correcta:




\begin{quote}
I used AI tools as coding assistants for scaffolding and iteration, while retaining responsibility for architecture, integration, debugging, validation and final implementation decisions.
\end{quote}




Forma incorrecta:




\begin{quote}
AI built the project.
\end{quote}




Forma igualmente incorrecta:




\begin{quote}
I did not use AI.
\end{quote}




\medskip\hrule\medskip




\subsection{13. Red flags que deben evitarse en documentos públicos}




\begin{center}
\scriptsize
\begin{longtable}{>{\raggedright\arraybackslash}p{0.455\linewidth} >{\raggedright\arraybackslash}p{0.455\linewidth}}
\toprule
Red flag & Por qué daña \\
\midrule
\endfirsthead
\toprule
Red flag & Por qué daña \\
\midrule
\endhead
Demasiados roles en headline & Parece perfil incoherente \\
Llenar CV con cursos & Sustituye evidencia por consumo de contenido \\
Presentar proyectos incompletos como productos & Riesgo en entrevista técnica \\
Usar métricas no verificadas & Riesgo reputacional \\
Publicar repos con claves, paths personales o código roto & Riesgo técnico/profesional \\
Decir “senior” por proyectos personales & Desalineación con reclutadores \\
Enfocar demasiado videojuegos sin portfolio de juegos & Reduce probabilidad \\
Hablar de EU como si ya hubiera permiso & Riesgo legal/credibilidad \\
Mezclar FitApp, ARA, TwinSight, VFX, full-stack y character art al mismo nivel & Dilución \\
No aclarar AI-assisted development si preguntan & Riesgo de percepción de engaño \\
\bottomrule
\end{longtable}
\end{center}




\medskip\hrule\medskip




\subsection{14. Decisiones que quedan abiertas}




Antes de cerrar estrategia final, se deben responder o investigar estas preguntas:




\subsubsection{14.1 Perfil y disponibilidad}




\begin{enumerate}
\item ¿Fecha estimada de sustentación?
\item ¿Disponibilidad real semanal después de sustentación?
\item ¿Runway financiero actualizado?
\item ¿Ingreso mínimo mensual neto necesario?
\item ¿Hay RUT activo?
\item ¿Hay facturación electrónica habilitada?
\item ¿Hay contador o asesor tributario?
\item ¿Hay Payoneer, Wise, Deel, Remote u otro setup?
\end{enumerate}




\subsubsection{14.2 TwinSight}




\begin{enumerate}
\item ¿Cuál es la URL final del repo correcto?
\item ¿Existe demo WebGL pública?
\item ¿Corre en móvil actualmente?
\item ¿Cuáles son las métricas finales exactas?
\item ¿Qué partes de assets son propias, CAD del fabricante o reconstruidas?
\item ¿Qué licencia/nota legal conviene usar?
\item ¿Qué screenshots ya existen?
\item ¿Hay video demo o debe producirse?
\item ¿Qué código fue IA-assisted y cómo explicarlo?
\end{enumerate}




\subsubsection{14.3 Proyectos secundarios}




\begin{enumerate}
\item ¿ARA corre end-to-end?
\item ¿Tiene API keys o datos sensibles que limpiar?
\item ¿ai-news-aggregator corre actualmente?
\item ¿FitApp tiene demo o solo planificación?
\item ¿El Blender portrait está terminado y renderizable?
\item ¿Qué repos están públicos ahora?
\item ¿Qué repos están pineados?
\end{enumerate}




\subsubsection{14.4 Cursos}




\begin{enumerate}
\item ¿Qué cursos están realmente completados?
\item ¿Qué cursos están en progreso?
\item ¿Qué cursos tienen output visible?
\item ¿Qué acceso Rebelway sigue disponible?
\item ¿Qué curso podría producir el portfolio artifact de mayor ROI en 2–4 semanas?
\item ¿Hay certificados emitidos por alguno? Si no, no se tratan como certificaciones.
\end{enumerate}




\subsubsection{14.5 Mercado}




\begin{enumerate}
\item ¿Qué roles tienen mayor frecuencia remota real desde LATAM?
\item ¿Qué títulos usan empresas de technical visualization?
\item ¿Qué roles pagan USD 3k/mes a perfiles early-career contractors?
\item ¿Qué empresas aceptan Colombia/LATAM contractor?
\item ¿Dónde encaja mejor TwinSight: games, XR, simulation, digital twin, product visualization, defense-adjacent, education, drones?
\end{enumerate}




Estas preguntas se resolverán en archivos posteriores, no en esta sección.




\medskip\hrule\medskip




\subsection{15. Implicaciones para archivos posteriores}




\begin{center}
\scriptsize
\begin{longtable}{>{\raggedright\arraybackslash}p{0.455\linewidth} >{\raggedright\arraybackslash}p{0.455\linewidth}}
\toprule
Archivo futuro & Regla derivada de esta fuente de verdad \\
\midrule
\endfirsthead
\toprule
Archivo futuro & Regla derivada de esta fuente de verdad \\
\midrule
\endhead
\texttt{02\_market\_role\_fit\_and\_positioning.md} & Validar roles contra mercado; no asumir que Technical Artist puro es el mejor sin datos \\
\texttt{03\_salary\_benchmark\_and\_remote\_colombia.md} & Separar employee/contractor, local/global, junior/mid; no usar rangos globales como si fueran remotos desde Colombia \\
\texttt{04\_eu\_portugal\_germany\_mobility\_strategy.md} & No afirmar derechos UE actuales; usar fuentes oficiales \\
\texttt{05\_language\_strategy\_and\_roi.md} & Inglés primero; alemán/portugués según horizonte; no sacrificar portfolio por idiomas en corto plazo \\
\texttt{06\_education\_courses\_rebelway\_masters\_phd.md} & Clasificar cursos por output/ROI; no por acumulación \\
\texttt{07\_portfolio\_strategy\_and\_project\_architecture.md} & TwinSight primero; ARA secundario; limitar dispersión \\
\texttt{08\_twinsight\_x500\_case\_study.md} & Basarse solo en hechos, métricas y assets verificables \\
\texttt{09\_ara\_ai\_tools\_and\_secondary\_projects.md} & Decidir inclusión por evidencia, no por potencial abstracto \\
\texttt{10\_github\_linkedin\_cv\_strategy.md} & Redactar desde claims defendibles; corregir tesis/ARA inconsistency \\
\texttt{11\_company\_targets\_job\_boards\_recruiters.md} & Buscar empresas que encajen con TwinSight y remoto Colombia, no solo nombres grandes \\
\texttt{12\_application\_system\_and\_tracker.md} & Sistema simple, medible, no sobreplanificado \\
\texttt{13\_interview\_negotiation\_and\_offer\_risk.md} & Preparar respuestas para low formal experience, AI-assisted work, salary, remote Colombia \\
\texttt{14\_30\_60\_90\_execution\_plan.md} & Plan desde estado actual, no desde cronogramas viejos \\
\texttt{15\_final\_integrated\_strategy.md} & Sintetizar lo validado; no introducir claims nuevos \\
\bottomrule
\end{longtable}
\end{center}




\medskip\hrule\medskip




\subsection{16. Criterios de aprobación de esta sección}




Esta sección se considera aprobada si:




\begin{enumerate}
\item fija que la tesis está completa y pendiente de sustentación;
\item separa hechos, supuestos, aspiraciones y claims prohibidos;
\item identifica TwinSight como activo principal;
\item mantiene ARA como secundario y no como eje principal;
\item trata cursos como aprendizaje, no certificación;
\item evita afirmar autorización laboral UE;
\item conserva rutas Portugal/Alemania como mediano-largo plazo;
\item limita la dispersión de proyectos;
\item define qué debe verificarse antes de CV, LinkedIn, GitHub y portfolio;
\item deja claro que el siguiente paso requiere investigación externa para roles/mercado.
\end{enumerate}




\medskip\hrule\medskip




\subsection{17. Próximo archivo recomendado}




El siguiente archivo debe ser:




\begin{mdcode}
Bloque de codigo: text
02_market_role_fit_and_positioning.md
\end{mdcode}




\subsubsection{Recomendación de herramientas para la siguiente sección}




\begin{center}
\scriptsize
\begin{longtable}{>{\raggedright\arraybackslash}p{0.305\linewidth} >{\raggedright\arraybackslash}p{0.305\linewidth} >{\raggedright\arraybackslash}p{0.305\linewidth}}
\toprule
Herramienta & Activar & Motivo \\
\midrule
\endfirsthead
\toprule
Herramienta & Activar & Motivo \\
\midrule
\endhead
Deep Research & Sí & Se necesita validar demanda actual, roles reales, títulos, keywords y fit por mercado. \\
Modo agente & Sí, si está disponible & Útil para recolectar ofertas reales, títulos frecuentes, requirements, remote restrictions y keywords desde job boards/careers pages. \\
Modo normal & Sí & Necesario para sintetizar y decidir jerarquía estratégica después de recolectar datos. \\
\bottomrule
\end{longtable}
\end{center}




\subsubsection{Objetivo de la siguiente sección}




Determinar, con evidencia actual, cuáles roles debe priorizar Alexander y cuáles deben quedar como secundarios o contingencia:




\begin{itemize}
\item Unity Technical Artist;
\item Real-Time 3D Developer;
\item Unity WebGL Developer;
\item Interactive 3D Developer;
\item Technical Visualization Developer;
\item CAD-to-Realtime Technical Artist;
\item XR / Simulation / Digital Twin Developer;
\item Game Technical Artist;
\item Tools Developer;
\item Python AI Tools Developer;
\item Full-stack secondary route.
\end{itemize}




La sección 02 debe producir una matriz de roles basada en:




\begin{itemize}
\item probabilidad de entrada;
\item salario potencial;
\item ajuste con TwinSight;
\item ajuste con habilidades reales;
\item remoto desde Colombia;
\item velocidad a primera entrevista;
\item riesgo de sobreventa;
\item evidencia de portfolio requerida;
\item keywords ATS;
\item mercados objetivo.
\end{itemize}



